\section{Endogenous Preferences and Portfolio Choice}
\label{sec:preferences}
\subsection{A Stopped Operating Economy}
We first state the continuous-state diffusion benchmark, including its boundary contract and control margins. The executed regional decision experiments use its separately specified, eight-date lottery-settlement counterpart in Section \ref{sec:decision_target}. A certificate for that counterpart is an economic statement about its declared transition protocol; transferring it to this diffusion additionally requires the target-error account developed below.
We now specify an economy in which the preference-adjustment margin is active in the transition law. Wealth is $X$, risk aversion is $u$, and the controls are consumption $c$, preference adjustment $\theta$, and risky-asset share $\pi$. Inside
$D=(1.2,2.8)\times(0.5,2.0)$, the state evolves as
\begin{align}
 du_t&=\theta_t\,dt+\sigma_u\,dZ_t^u,\label{eq:ndu_u}\\
 dX_t&=\{(r+\pi_t(\mu-r))X_t-c_t\}\,dt
           +\pi_tX_t\sigma_X\,dZ_t^X,\qquad
 d\langle Z^u,Z^X\rangle_t=\varrho\,dt.
 \label{eq:ndu_x}
\end{align}
The operating contract terminates at $\tau=T\wedge\inf\{t:S_t\notin D\}$. At that time the agent receives the cardinal payoff
\begin{equation}
 G(u,X)=-0.02(u-2)^2+0.1\log X.
 \label{eq:ndu_terminal}
\end{equation}
The payoff summarizes liquidation or the outside option after leaving the operating regime; it is not a transfer restoring the state to the interior.

The feasible box is $c\in[0.05,0.8]$, $\theta\in[-0.2,0.2]$, and $\pi\in[-0.5,0.8]$. We set
\begin{equation}
 T=1,\quad \rho=0.04,\quad r=0.02,\quad \mu=0.08,
 \quad \sigma_X=0.2,\quad \sigma_u=0.05,\quad \varrho=-0.25.
 \label{eq:ndu_parameters}
\end{equation}
The baseline adjustment-cost coefficient is $k=2$. For a feasible policy,
\begin{equation}
 J_k^\pi(t,u,X)=\E^\pi_{t,u,X}\left[
 \int_t^\tau e^{-\rho(s-t)}
 \left\{\frac{c_s^{1-u_s}}{1-u_s}-\frac{k}{2}\theta_s^2\right\}ds
 +e^{-\rho(\tau-t)}G(u_\tau,X_\tau)\right].
 \label{eq:ndu_objective}
\end{equation}
Consumption is measured in the fixed units appearing in \eqref{eq:ndu_objective}. An additive or multiplicative normalization depending on $u$ changes incentives once $u$ is controlled. We therefore do not treat such transformations as innocuous changes of utility representation.

This boundary contract is a substantive specification. Constant preference volatility cannot remain in a closed interval without an exit or regulator mechanism. Moreover, under the displayed consumption bounds, the drift at $X=0.5$ is negative even at the largest feasible portfolio share: its maximum is $-0.016$. An unregulated state-constrained interpretation is consequently unavailable. Stopping repairs that problem without modifying the interior diffusion, imposing an unreported reflection, or injecting wealth. It is not equivalent to the earlier clipped transition.

The stopped problem has bounded payoffs and Lipschitz coefficients before exit. A useful check is that $G$ is a strict supersolution for every feasible control. The bounds in \eqref{eq:ndu_parameters} imply
\begin{equation}
 \frac{c^{1-u}}{1-u}-\frac{k\theta^2}{2}
 +\cL^{c,\theta,\pi}G-\rho G\leq -0.8162
 \qquad(k\geq0).
 \label{eq:superG}
\end{equation}
The supplement derives the bound without a grid maximization. Together with a fixed nondegenerate admissible diffusion, it supplies a boundary-continuity argument for the stopped value. The comparison also has economic content: continued operation has a cost relative to the liquidation payoff. Exit incentives are therefore part of this calibration and must be reported, rather than interpreted as a generic wealth or risk-aversion effect.

\subsection{The Operator and Control Margins}
Writing subscripts for derivatives, the HJB equation in the interior is
\begin{align}
 -V_t=\sup_{c,\theta,\pi}\big\{
 &\tfrac{c^{1-u}}{1-u}-\tfrac12k\theta^2-\rho V
 +\theta V_u+[(r+\pi(\mu-r))X-c]V_X\notag\\
 &+\tfrac12\sigma_u^2V_{uu}
 +\tfrac12\pi^2X^2\sigma_X^2V_{XX}
 +\varrho\sigma_u\sigma_X\pi X V_{uX}\big\},
 \label{eq:ndu_HJB}
\end{align}
with $V=G$ on the spatial boundary and at $T$. Both diagonal diffusion terms and the mixed term belong to the operator. In particular, the generator of $u^2$ contains $\sigma_u^2$ even if the preference drift is zero.

For an interior optimum, the first-order conditions are
\begin{align}
 c^{-u}&=V_X,& k\theta&=V_u,&
 0&=(\mu-r)XV_X+\pi X^2\sigma_X^2V_{XX}
               +\varrho\sigma_u\sigma_X X V_{uX}.
 \label{eq:ndu_FOC}
\end{align}
When $V_X>0$, $V_{XX}<0$, and all relevant controls are interior, these give
\begin{equation}
 c=V_X^{-1/u},\qquad \theta=V_u/k,\qquad
 \pi=-\frac{(\mu-r)V_X}{X\sigma_X^2V_{XX}}
      -\frac{\varrho\sigma_uV_{uX}}{X\sigma_XV_{XX}}.
 \label{eq:ndu_controls}
\end{equation}
The final expression separates a mean--variance term from a preference-hedging term. Its sign is not determined by $\varrho<0$ alone: the cross derivative and curvature also matter. Nor need the stopped economy satisfy $V_X>0$ throughout the domain. The constrained implementation compares feasible controls rather than mechanically applying these formulas where their premises fail. At a binding box constraint, the appropriate one-sided derivative or complementary-slackness condition replaces an interior equality.

The cardinal utility derivative is
\begin{equation}
 \partial_u\left(\frac{c^{1-u}}{1-u}\right)
 =\frac{c^{1-u}\{1-(1-u)\log c\}}{(1-u)^2}.
 \label{eq:normalization_derivative}
\end{equation}
It need not have the sign suggested by a verbal interpretation of increased risk aversion. The full shadow value $V_u$ incorporates this derivative, future opportunities, adjustment costs, covariance, and liquidation. The numerical policy is informative only relative to these declared primitives.

\subsection{A Comparative Static That Survives Endogenous Policies}
Fix an initial state and let
\begin{equation}
 C^\pi=\E^\pi\int_0^\tau e^{-\rho t}\frac{\theta_t^2}{2}\,dt,
 \qquad B^\pi=J_k^\pi+kC^\pi.
 \label{eq:effort_functional}
\end{equation}
For a fixed policy, $B^\pi$ and $C^\pi$ are independent of $k$: changing the coefficient does not change the state law unless the policy is changed. In particular, the stopping time is allowed to be endogenous to the policy.

\begin{theorem}[Adjustment costs and cumulative effort]
\label{thm:k}
Suppose the feasible policy class and all primitives other than $k$ are fixed, and optimizers exist at the coefficients being compared. Then
$V(k)=\sup_\pi\{B^\pi-kC^\pi\}$ is decreasing and convex. If $k_2>k_1$ and $\pi_i$ is optimal at $k_i$, then
\begin{equation}
 C^{\pi_2}\leq C^{\pi_1}.
 \label{eq:effort_monotonicity}
\end{equation}
Each $-C^{\pi_i}$ is a subgradient of $V$ at $k_i$. At differentiability points, all optimal effort levels equal $-V'(k)$. For policies with objective losses at most $e_1,e_2$ at the respective coefficients,
\begin{equation}
 C^{\pi_2}-C^{\pi_1}\leq\frac{e_1+e_2}{k_2-k_1}.
 \label{eq:effort_approximate}
\end{equation}
\end{theorem}

The proof uses the two optimality inequalities and is given in Supplement S.3. The approximate version is useful when numerical policies are used in an economic comparison: the effort ordering can be assessed against a welfare-error budget. The result does not hold pointwise for $\theta(t,u,X)$ merely because \eqref{eq:ndu_controls} contains $1/k$. The shadow value, portfolio choice, consumption, and exit distribution all change when the problem is re-solved.

The theorem is an economic comparative static, not an identification result. Inferring $k$ from data additionally requires an observation model for latent preferences, an identified moment map, and control of numerical error in that map. The solver supplies policy and value objects for such a model; a graph of the policy alone cannot establish identification. This distinction leaves the structural-estimation motivation intact while identifying the additional objects it would require.
